% Sections 2--3: formal setup and theoretical analysis.
% Required packages: amsmath, amssymb, amsthm, bm, mathtools.
% Suggested theorem environments:
% \newtheorem{definition}{Definition}
% \newtheorem{theorem}{Theorem}
% \newtheorem{proposition}{Proposition}
% \newtheorem{corollary}{Corollary}

\section{Problem Formulation}
\label{sec:formulation}

\subsection{Longitudinal treatment setting}

We consider a patient observed at discrete decision times
$t=0,\ldots,T$. Let $O_t\in\mathcal O$ denote the observed clinical
state, $A_t\in\mathcal A$ the treatment administered at time $t$, and
$Y_t=q(O_t)\in\mathbb R$ an outcome of interest. The history available
immediately before treatment $A_t$ is
\begin{equation}
    H_t=(O_0,A_0,O_1,\ldots,A_{t-1},O_t)\in\mathcal H.
    \label{eq:history}
\end{equation}
The observational data are generated under a behaviour policy $\mu$,
such that $A_t\sim\mu(\cdot\mid H_t)$. Our training set contains $N$
factual patient trajectories,
\begin{equation}
    \mathcal D_\mu
    =\left\{
    (O^{(i)}_0,A^{(i)}_0,\ldots,A^{(i)}_{T_i-1},O^{(i)}_{T_i})
    \right\}_{i=1}^{N}.
    \label{eq:dataset}
\end{equation}

For a prediction horizon $H$, write
$\bm a_{t:t+H-1}=(a_t,\ldots,a_{t+H-1})\in\mathcal A^H$ for an
externally specified treatment schedule. The potential outcome at
horizon $h\leq H$ is denoted by
\begin{equation}
    Y_{t+h}(\bm a_{t:t+h-1}).
    \label{eq:potential-outcome}
\end{equation}
For two schedules $\bm a$ and $\bm a'$, the conditional longitudinal
treatment-effect trajectory is
\begin{align}
    \tau_h(H_t;\bm a,\bm a')
    &:=
    \mathbb E\!\left[
        Y_{t+h}(\bm a_{t:t+h-1})
        -Y_{t+h}(\bm a'_{t:t+h-1})
        \mid H_t
    \right],
    \label{eq:true-effect}\\
    \bm\tau_{1:H}(H_t;\bm a,\bm a')
    &:= (\tau_1,\ldots,\tau_H).
    \label{eq:true-effect-trajectory}
\end{align}
Unlike a scalar terminal effect or policy value,
$\bm\tau_{1:H}$ retains when an effect emerges, accumulates, persists,
or decays.

The causal interpretation of \eqref{eq:true-effect} requires the usual
longitudinal identification assumptions: consistency, sequential
exchangeability given the measured history, and positivity for the
state--action relationships under consideration. These assumptions hold
by construction in our controlled simulator. Our representation-learning
method does not, by itself, establish them in observational clinical data.

\subsection{Medical world models as longitudinal simulators}

A latent medical world model consists of an encoder $E_\theta$, an
action-conditioned transition $F_\theta$, and an observation decoder
$G_\theta$:
\begin{align}
    Z_t &= E_\theta(H_t),
    \label{eq:encoder}\\
    \widehat Z_{t+1} &= F_\theta(\widehat Z_t,A_t),
    \label{eq:transition}\\
    \widehat O_{t+1} &= G_\theta(\widehat Z_{t+1}).
    \label{eq:decoder}
\end{align}
Starting from $\widehat Z_t=E_\theta(H_t)$, the model is recursively
rolled forward under a candidate schedule:
\begin{equation}
    \widehat Z_{t+h}(\bm a)
    =F_\theta^{(h)}
    \!\left(E_\theta(H_t),\bm a_{t:t+h-1}\right),
    \qquad
    \widehat Y_{t+h}(\bm a)
    =q\!\left(G_\theta(\widehat Z_{t+h}(\bm a))\right).
    \label{eq:model-rollout}
\end{equation}
The world-model estimate of the treatment-effect trajectory is obtained
by contrasting two rollouts from the same encoded history:
\begin{equation}
    \widehat\tau_h(H_t;\bm a,\bm a')
    =\widehat Y_{t+h}(\bm a)-\widehat Y_{t+h}(\bm a').
    \label{eq:model-effect}
\end{equation}

We evaluate intervention fidelity using the trajectory-level effect risk
\begin{equation}
    \mathcal R_{\mathrm{eff}}^{(H)}(\theta;\Pi)
    =
    \mathbb E_{H_t,(\bm a,\bm a')\sim\Pi}
    \left[
        \frac{1}{H}\sum_{h=1}^{H}
        \bigl(\widehat\tau_h-\tau_h\bigr)^2
    \right],
    \label{eq:effect-risk}
\end{equation}
where $\Pi$ is an evaluation distribution over pairs of schedules. In
the controlled experiments, both potential-outcome trajectories are
generated by branching the same simulator state, so
\eqref{eq:effect-risk} is exactly evaluable. Counterfactual outcomes are
used for evaluation only and are never supplied to the proposed training
objective.

\subsection{Factual predictive objectives}

Conventional one-step training minimises factual prediction risk under
the behaviour policy:
\begin{equation}
    \mathcal L_{\mathrm{1step}}(\theta)
    =\mathbb E_{\mathcal D_\mu}
    \left[
        \ell\!\left(
            G_\theta(F_\theta(E_\theta(H_t),A_t)),O_{t+1}
        \right)
    \right].
    \label{eq:one-step-loss}
\end{equation}
A stronger closed-loop objective rolls the model forward under the
observed action sequence and penalises multiple future observations:
\begin{equation}
    \mathcal L_{\mathrm{roll}}^{(H)}(\theta)
    =\mathbb E_{\mathcal D_\mu}
    \left[
        \sum_{h=1}^{H}w_h
        \ell\!\left(
            G_\theta(\widehat Z_{t+h}(\bm A^{\mathrm{obs}})),
            O_{t+h}
        \right)
    \right].
    \label{eq:rollout-loss}
\end{equation}
This objective directly addresses multi-step observable drift, but it
still evaluates only factual schedules drawn from $\mu$ and constrains
the latent state only through the decoder. Our theoretical question is
therefore not whether perfect prediction of both interventional
trajectories would imply a correct effect---it would---but whether low
factual predictive risk is sufficient to certify a recursively composed
latent simulator under alternative schedules.

\section{When Predictive Fidelity Is Insufficient}
\label{sec:theory}

We isolate three distinct obstacles: representation aliasing, error
propagation through recursive transitions, and treatment-policy shift.
The results are diagnostic rather than identification theorems: they
characterise why a statistically identified longitudinal effect may
still be poorly realised by an imperfect latent world model.

\subsection{Representation aliasing}

\begin{definition}[Treatment-sufficient representation]
\label{def:treatment-sufficient}
Let $\Pi_H\subseteq\mathcal A^H$ be a supported class of schedules. An
encoder $E:\mathcal H\to\mathcal Z$ is treatment-sufficient up to
horizon $H$ if, for any histories $H_t,H_t'$,
\begin{equation}
    E(H_t)=E(H_t')
    \quad\Longrightarrow\quad
    \mathbb E[Y_{t+h}(\bm a)\mid H_t]
    =
    \mathbb E[Y_{t+h}(\bm a)\mid H_t']
    \label{eq:treatment-sufficiency}
\end{equation}
for every $\bm a\in\Pi_H$ and every $h\leq H$.
\end{definition}

Treatment sufficiency is stronger than sufficiency for next-step factual
prediction: two histories may have similar futures under the behaviour
policy while responding differently to an alternative treatment.

\begin{theorem}[Irreducible effect error under latent aliasing]
\label{thm:aliasing}
Suppose that two histories $H$ and $H'$ are aliased by the encoder,
$E(H)=E(H')=z$, but for some schedule pair $(\bm a,\bm a')$ and horizon
$h$ their true effects differ by
\begin{equation}
    \Delta
    =\left|
        \tau_h(H;\bm a,\bm a')
        -\tau_h(H';\bm a,\bm a')
    \right|>0.
    \label{eq:effect-separation}
\end{equation}
Then every effect estimator that is a function only of $z,\bm a,\bm a'$
satisfies
\begin{equation}
    \max\!\left{
      |\widehat\tau_h(z)-\tau_h(H)|,
      |\widehat\tau_h(z)-\tau_h(H')|
    \right}
    \geq \frac{\Delta}{2}.
    \label{eq:aliasing-max-bound}
\end{equation}
If the two histories have equal probability, their average squared
effect error is at least $\Delta^2/4$.
\end{theorem}

\begin{proof}
Because the estimator receives the same latent state for both histories,
it must output a common value $c$. The triangle inequality gives
\begin{equation}
    \Delta
    =|\tau_h(H)-\tau_h(H')|
    \leq |\tau_h(H)-c|+|c-\tau_h(H')|,
\end{equation}
so at least one term is no smaller than $\Delta/2$. Moreover,
\begin{equation}
    \frac{1}{2}\bigl[(c-u)^2+(c-v)^2\bigr]
    =\left(c-\frac{u+v}{2}\right)^2+rac{(u-v)^2}{4}
    \geq\frac{\Delta^2}{4},
\end{equation}
where $u=\tau_h(H)$ and $v=\tau_h(H')$.
\end{proof}

Theorem~\ref{thm:aliasing} identifies a representation-level failure:
once histories with different intervention responses are merged, no
downstream transition or decoder can recover the missing distinction.
Increasing latent dimension may reduce aliasing, but dimension alone
does not guarantee treatment sufficiency.

\subsection{Propagation through multi-step transitions}

We next separate information availability from its temporal use. Let
$Z^*_{t+h}(\bm a)$ denote a treatment-sufficient target state along the
true trajectory and suppose that
\begin{equation}
    Z^*_{t+k+1}(\bm a)
    =F^*(Z^*_{t+k}(\bm a),a_{t+k}).
    \label{eq:true-latent-transition}
\end{equation}
Define the rollout error and the local transition approximation error by
\begin{align}
    e_k(\bm a)
    &:=\|\widehat Z_{t+k}(\bm a)-Z^*_{t+k}(\bm a)\|,
    \label{eq:rollout-latent-error}\\
    \epsilon_k(\bm a)
    &:=\|F_\theta(Z^*_{t+k}(\bm a),a_{t+k})
             -Z^*_{t+k+1}(\bm a)\|.
    \label{eq:local-transition-error}
\end{align}

\begin{theorem}[Recursive latent-error propagation]
\label{thm:propagation}
Assume $F_\theta$ is $L_F$-Lipschitz in its latent-state argument:
\begin{equation}
    \|F_\theta(z,a)-F_\theta(z',a)\|
    \leq L_F\|z-z'\|.
    \label{eq:transition-lipschitz}
\end{equation}
Then, for every treatment schedule $\bm a$,
\begin{align}
    e_{k+1}(\bm a)
    &\leq L_F e_k(\bm a)+\epsilon_k(\bm a),
    \label{eq:one-step-recursion}\\
    e_h(\bm a)
    &\leq L_F^h e_0
      +\sum_{k=0}^{h-1}L_F^{h-1-k}\epsilon_k(\bm a).
    \label{eq:multi-step-recursion}
\end{align}
\end{theorem}

\begin{proof}
Add and subtract $F_\theta(Z^*_{t+k},a_{t+k})$ and apply the triangle
inequality and \eqref{eq:transition-lipschitz}:
\begin{align}
e_{k+1}
&=\|F_\theta(\widehat Z_{t+k},a_{t+k})-Z^*_{t+k+1}\|\\
&\leq
\|F_\theta(\widehat Z_{t+k},a_{t+k})
  -F_\theta(Z^*_{t+k},a_{t+k})\|
+\epsilon_k\\
&\leq L_F e_k+\epsilon_k.
\end{align}
Recursively expanding this inequality yields
\eqref{eq:multi-step-recursion}.
\end{proof}

The bound distinguishes stable ($L_F<1$), additive ($L_F=1$), and
potentially amplifying ($L_F>1$) rollout regimes. Even if the initial
state contains the required information, small local transition errors
can therefore accumulate or amplify across treatment decisions.

\begin{corollary}[Longitudinal effect-error bound]
\label{cor:effect-bound}
Let $g=q\circ G_\theta$ be $L_G$-Lipschitz. Suppose its approximation
error on the treatment-sufficient state is bounded by
\begin{equation}
    |g(Z^*_{t+h}(\bm a))-Y_{t+h}(\bm a)|
    \leq\eta_h(\bm a).
    \label{eq:decoder-residual}
\end{equation}
Then the effect error for two schedules is bounded by
\begin{align}
|\widehat\tau_h-\tau_h|
&\leq
L_G\bigl[e_h(\bm a)+e_h(\bm a')\bigr]
+\eta_h(\bm a)+\eta_h(\bm a')
\label{eq:effect-latent-bound}\\
&\leq
L_G\!\sum_{k=0}^{h-1}L_F^{h-1-k}
\bigl[\epsilon_k(\bm a)+\epsilon_k(\bm a')\bigr]
+2L_GL_F^h e_0
+\eta_h(\bm a)+\eta_h(\bm a').
\label{eq:effect-transition-bound}
\end{align}
\end{corollary}

\begin{proof}
Apply the triangle inequality separately to the two potential-outcome
predictions, use the Lipschitz condition on $g$, and substitute
Theorem~\ref{thm:propagation}.
\end{proof}

Corollary~\ref{cor:effect-bound} formalises why accurate local response
mapping is not sufficient for an accurate longitudinal effect curve:
both treatment branches depend on recursively propagated latent states.

\subsection{Treatment-policy shift}

Let $\rho_\mu(s,a)$ and $\rho_\pi(s,a)$ denote finite-horizon
state--action occupancy distributions under the behaviour policy and a
target policy, respectively. For a non-negative local transition loss
$\ell_\theta(s,a)$, factual training controls
\begin{equation}
    \mathcal R_\mu(\theta)
    =\mathbb E_{(s,a)\sim\rho_\mu}[\ell_\theta(s,a)],
    \label{eq:behavior-risk}
\end{equation}
whereas deployment depends on
\begin{equation}
    \mathcal R_\pi(\theta)
    =\mathbb E_{(s,a)\sim\rho_\pi}[\ell_\theta(s,a)].
    \label{eq:target-risk}
\end{equation}

\begin{proposition}[Occupancy-shift bound]
\label{prop:shift}
If $\rho_\pi$ is absolutely continuous with respect to $\rho_\mu$ and
$w(s,a)=\rho_\pi(s,a)/\rho_\mu(s,a)$, then
\begin{equation}
    \mathcal R_\pi(\theta)
    =\mathbb E_{\rho_\mu}[w\ell_\theta]
    \leq
    \sqrt{\mathbb E_{\rho_\mu}[w^2]}
    \sqrt{\mathbb E_{\rho_\mu}[\ell_\theta^2]}.
    \label{eq:shift-bound}
\end{equation}
\end{proposition}

\begin{proof}
The equality follows by change of measure and the inequality follows
from Cauchy--Schwarz.
\end{proof}

Thus low factual transition risk need not imply low risk under a target
treatment policy when the occupancy ratio is highly variable. If
$\rho_\pi(s,a)>0$ where $\rho_\mu(s,a)=0$, no such guarantee is
available without additional structural assumptions. Our compositional
evaluation therefore distinguishes unsupported extrapolation from the
more informative case in which local state--action relationships are
covered but the complete temporal schedule is held out.

\section{Action-Conditioned Temporal JEPA}
\label{sec:method}

The analysis in Section~\ref{sec:theory} identifies two model-side
requirements for reliable longitudinal effect estimation. First, the
encoded state must preserve distinctions that change future treatment
responses. Second, these distinctions must remain coherent when the
transition is recursively composed. The first requirement is a property
of the learned representation and available data; here we target the
second by directly supervising multi-horizon latent evolution.

\subsection{Motivation}

Observable rollout loss constrains $G_\theta(\widehat Z_{t+h})$, but a
many-to-one decoder can map distinct latent states to similar observed
outcomes. It therefore need not directly minimise the latent errors
$e_h$ appearing in Corollary~\ref{cor:effect-bound}. We introduce an
action-conditioned multi-horizon latent objective that complements
observable rollout training.

\subsection{Multi-horizon latent objective}

For a factual training window, the online encoder first represents the
observed history $H_t$. The transition model then recursively predicts
future latent states using the treatments that were actually observed in
that window. A slowly updated target encoder represents the corresponding
future histories:
\begin{align}
    \widetilde Z_{t+h}
    &=F_\theta^{(h)}
      \!\left(E_\theta(H_t),A_{t:t+h-1}^{\mathrm{obs}}\right),
    \label{eq:jepa-prediction}\\
    Z^*_{t+h}
    &=E_{\bar\theta}(H_{t+h}),
    \label{eq:jepa-target}\\
    \mathcal L_{\mathrm{TJEPA}}
    &=\mathbb E_{\mathcal D_\mu}
      \left[
      \sum_{h\in\mathcal K}
      d_{\mathcal Z}\!\left(
          \widetilde Z_{t+h},
          \operatorname{sg}(Z^*_{t+h})
      \right)
      \right],
    \qquad \mathcal K=\{1,2,4,8\},
    \label{eq:jepa-loss}
\end{align}
where $E_{\bar\theta}$ is an exponential-moving-average target encoder
and $\operatorname{sg}$ denotes stop-gradient. The complete objective is
\begin{equation}
    \mathcal L(\theta)
    =\mathcal L_{\mathrm{roll}}^{(H)}(\theta)
     +\lambda_J\mathcal L_{\mathrm{TJEPA}}(\theta).
    \label{eq:complete-objective}
\end{equation}

No simulator state, alternative-treatment outcome, treatment-effect
label, or policy-value target is used in
\eqref{eq:jepa-loss}. At deployment, the target encoder is discarded;
the learned encoder, transition, and decoder are rolled forward under an
externally specified candidate treatment schedule.

\subsection{Scope of the method}

Temporal JEPA does not create counterfactual information or solve causal
identification. Its narrower role is to use available factual future
histories to constrain the multi-step latent propagation term that
appears in \eqref{eq:effect-latent-bound}. Consequently, a reduction in
$\mathcal L_{\mathrm{TJEPA}}$ is informative for longitudinal effects
only when the target representation is treatment-sufficient, the
decoder retains the relevant outcome information, and the evaluated
treatment schedules have adequate support. These conditions motivate
our separate experiments on latent capacity, rollout horizon, and
treatment-policy shift.
